% sec2_2.tex — Section 2.2: Fundamental Concepts in Time-Series Analysis

In this section, we define the very basic concepts in time-series analysis that will
form an integral part of our language.  The key concept is a \textbf{stochastic process},
which is just a fancy name for a sequence of random variables.  If the index for
the random variables is interpreted as representing time, the stochastic process is
called a \textbf{time series}.  If $\{z_i\}$ $(i = 1, 2, \ldots)$ is a stochastic process,
its \textbf{realization} or a \textbf{sample path} is an assignment to each $i$ of a
possible value of $z_i$.  So a realization of $\{z_i\}$ is a sequence of real numbers.
We will frequently use the term time series to mean both the realization and the
process of which it is a realization.

%% ─── Need for Ergodic Stationarity ────────────────────────────────────────
\subsection*{Need for Ergodic Stationarity}

The fundamental problem in time-series analysis is that we can observe the
realization of the process only once.  For example, the sample on the U.S.\ annual
inflation rate for the period from 1946 to 1995 is a string of 50 particular numbers,
which is just one possible outcome of the underlying stochastic process for the
inflation rate; if history took a different course, we could have obtained a different
sample.

If we could observe the history many times over, we could assemble many samples,
each containing a possibly different string of 50 numbers.  The mean inflation rate
for, say, 1995 can then be estimated by taking the average of the 1995 inflation
rate (the 50th element of the string) across those samples.  The population mean
thus estimated is sometimes called the \textbf{ensemble mean}.  In the language of
general equilibrium theory in economics, the ensemble mean is the average across
all the possible different states of nature at any given calendar time.

Of course, it is not feasible to observe many different alternative histories.  But
if the distribution of the inflation rate remains unchanged (this property will be
referred to as \textbf{stationarity}), the particular string of 50 numbers we do observe
can be viewed as 50 different values from the \emph{same} distribution.  Furthermore,
if the process is not too persistent (what's called \textbf{ergodicity} has this property),
each element of the string will contain some information not available from the
other elements, and, as shown below, the time average over the elements of the
single string will be consistent for the ensemble mean.

%% ─── Various Classes of Stochastic Processes ──────────────────────────────
\subsection*{Various Classes of Stochastic Processes\protect\footnote{Many of the
concepts collected in this subsection can also be found in Section~4.7 of Davidson
and MacKinnon (1993).}}

\paragraph{Stationary Processes.}
A stochastic process $\{\mathbf{z}_i\}$ $(i = 1, 2, \ldots)$ is (strictly)
\textbf{stationary} if, for any given finite integer $r$ and for any set of subscripts,
$i_1, i_2, \ldots, i_r$, the joint distribution of
$(\mathbf{z}_{i_1}, \mathbf{z}_{i_2}, \ldots, \mathbf{z}_{i_r})$ depends only on
$i_1 - i$, $i_2 - i$, $i_3 - i, \ldots, i_r - i$ but not on $i$.  For example, the
joint distribution of $(\mathbf{z}_1, \mathbf{z}_5)$ is the same as that of
$(\mathbf{z}_{12}, \mathbf{z}_{16})$.  What matters for the distribution is the
relative position in the sequence.  In particular, the distribution of $\mathbf{z}_i$
does not depend on the absolute position, $i$, of $\mathbf{z}_i$, so the mean,
variance, and other higher moments, if they exist, remain the same across $i$.
The definition also implies that any transformation (function) of a stationary
process is itself stationary, that is, if $\{\mathbf{z}_i\}$ is stationary, then
$\{f(\mathbf{z}_i)\}$ is.\footnote{The function $f(\cdot)$ needs to be ``measurable''
so that $f(\mathbf{z}_i)$ is a well-defined random variable.  Any continuous function
is measurable.  In what follows, we won't bother to add the qualifier ``measurable''
when a function $f$ of a random variable is understood to be a random variable.}
For example, $\{\mathbf{z}_i \mathbf{z}_i'\}$ is stationary if $\{\mathbf{z}_i\}$ is.

\begin{example}[i.i.d.\ sequences]
\label{ex:2.1}
A sequence of independent and identically distributed random variables is a
stationary process that exhibits no serial dependence.
\end{example}

\begin{example}[constant series]
\label{ex:2.2}
Draw $z_1$ from some distribution and then set $z_i = z_1$
$(i = 2, 3, \ldots)$.  So the value of the process is frozen at the initial date.
The process $\{z_i\}$ thus created is a stationary process that exhibits maximum
serial dependence.
\end{example}

Evidently, if a vector process $\{\mathbf{z}_i\}$ is stationary, then each element of
the vector forms a univariate stationary process.  The converse, however, is not true.

\begin{example}[element-wise vs.\ joint stationarity]
\label{ex:2.3}
Let $\{\varepsilon_i\}$ $(i = 1, 2, \ldots)$ be a scalar i.i.d.\ process.  Create a
two-dimensional process $\{\mathbf{z}_i\}$ from it by defining $z_{i1} = \varepsilon_i$
and $z_{i2} = \varepsilon_1$.  The scalar process $\{z_{i1}\}$ is stationary (this is
the process of Example~\ref{ex:2.1}).  The scalar process $\{z_{i2}\}$, too, is
stationary (the Example~\ref{ex:2.2} process).  The vector process
$\{\mathbf{z}_i\}$, however, is not (jointly) stationary, because the (joint)
distribution of $\mathbf{z}_1$ $(= (\varepsilon_1, \varepsilon_1)')$ differs from that
of $\mathbf{z}_2$ $(= (\varepsilon_2, \varepsilon_1)')$.
\end{example}

Most aggregate time series such as GDP are not stationary because they exhibit
time trends.  A less obvious example of nonstationarity is international exchange
rates, which are alleged to have increasing variance.  But many time series with
trend can be reduced to stationary processes.  A process is called
\textbf{trend stationary} if it is stationary after subtracting from it a (usually
linear) function of time (which is the index $i$).  If a process is not stationary
but its first difference, $z_i - z_{i-1}$, is stationary, $\{z_i\}$ is called
\textbf{difference stationary}.  Trend-stationary and difference-stationary processes
will be studied in Chapter~9.

\paragraph{Covariance Stationary Processes.}
A stochastic process $\{\mathbf{z}_i\}$ is \textbf{weakly} (or \textbf{covariance})
\textbf{stationary} if:
\begin{enumerate}[label=(\roman*)]
  \item $\E(\mathbf{z}_i)$ does not depend on $i$, and
  \item $\Cov(\mathbf{z}_i, \mathbf{z}_{i-j})$ exists, is finite, and depends only on
        $j$ but not on $i$ (for example, $\Cov(\mathbf{z}_1, \mathbf{z}_5)$ equals
        $\Cov(\mathbf{z}_{12}, \mathbf{z}_{16})$).
\end{enumerate}
The relative, not absolute, position in the sequence matters for the mean and
covariance of a covariance-stationary process.  Evidently, if a sequence is
(strictly) stationary \emph{and} if the variance and covariances are finite, then
the sequence is weakly stationary (hence the term ``strict'').  An example of a
covariance-stationary but not strictly stationary process will be given in
Example~\ref{ex:2.4} below.

The \textbf{$j$-th order autocovariance}, denoted $\bm{\Gamma}_j$, is defined as
\[
  \bm{\Gamma}_j \equiv \Cov(\mathbf{z}_i, \mathbf{z}_{i-j})
  \quad (j = 0, 1, 2, \ldots).
\]
The term ``auto'' comes about because the two random variables are taken from the
same process.  $\bm{\Gamma}_j$ does not depend on $i$ because of covariance
stationarity.  Also by covariance stationarity, $\bm{\Gamma}_j$ satisfies
\begin{equation}
  \bm{\Gamma}_j = \bm{\Gamma}'_{-j}.
  \label{eq:2.2.1}
\end{equation}
(Showing this is a review question below.)  The 0-th order autocovariance is the
variance $\bm{\Gamma}_0 = \Var(\mathbf{z}_i)$.  The processes in
Examples~\ref{ex:2.1} and \ref{ex:2.2} are covariance stationary if the variance of
the distribution is finite.  For the process of Example~\ref{ex:2.1},
$\bm{\Gamma}_0$ is the variance of the distribution and $\bm{\Gamma}_j = \mathbf{0}$
for $j \geq 1$.  For the process of Example~\ref{ex:2.2},
$\bm{\Gamma}_j = \bm{\Gamma}_0$.

For a scalar covariance stationary process $\{z_i\}$, the $j$-th order autocovariance
is now a scalar.  If $\gamma_j$ is this autocovariance, it satisfies
\begin{equation}
  \gamma_j = \gamma_{-j}.
  \label{eq:2.2.2}
\end{equation}

Take a string of $n$ successive values, $(z_i, z_{i+1}, \ldots, z_{i+n-1})$, from a
scalar covariance process.  By covariance stationarity, its $n \times n$
variance-covariance matrix is the same as that of $(z_1, z_2, \ldots, z_n)$ and is a
band spectrum matrix:
\[
  \Var(z_i, z_{i+1}, \ldots, z_{i+n-1})
  = \begin{bmatrix}
      \gamma_0    & \gamma_1    & \gamma_2 & \cdots & \gamma_{n-1} \\
      \gamma_1    & \gamma_0    & \gamma_1 & \cdots & \gamma_{n-2} \\
      \vdots      & \ddots      & \ddots   & \ddots & \vdots       \\
      \gamma_{n-2}& \cdots      & \gamma_1 & \gamma_0 & \gamma_1   \\
      \gamma_{n-1}& \cdots      & \gamma_2 & \gamma_1 & \gamma_0
    \end{bmatrix}.
\]
This is called the \textbf{autocovariance matrix} of the process.  The
\textbf{$j$-th order autocorrelation coefficient}, $\rho_j$, is defined as
\begin{equation}
  j\text{-th order autocorrelation coefficient}
  \equiv \rho_j
  \equiv \frac{\gamma_j}{\gamma_0}
  = \frac{\Cov(z_i, z_{i-j})}{\Var(z_i)}
  \quad (j = 1, 2, \ldots).
  \label{eq:2.2.3}
\end{equation}
For $j = 0$, $\rho_j = 1$.  The plot of $\{\rho_j\}$ against
$j = 0, 1, 2, \ldots$ is called the \textbf{correlogram}.

%% ─── White Noise Processes ─────────────────────────────────────────────────
\paragraph{White Noise Processes.}
A very important class of weakly stationary processes is a \textbf{white noise process},
a process with zero mean and no serial correlation:
\begin{gather*}
  \text{a covariance-stationary process $\{\mathbf{z}_i\}$ is \textbf{white noise} if}
  \\
  \E(\mathbf{z}_i) = \mathbf{0} \quad \text{and} \quad
  \Cov(\mathbf{z}_i, \mathbf{z}_{i-j}) = \mathbf{0} \;\text{for}\; j \neq 0.
\end{gather*}
Clearly, an independently and identically distributed (i.i.d.)\ sequence with mean
zero and finite variance is a special case of a white noise process.  For this reason,
it is called an \textbf{independent white noise process}.

\begin{example}[a white noise process that is not strictly
stationary\protect\footnote{Drawn from Example~7.8 of Anderson (1971, p.~379).}]
\label{ex:2.4}
Let $w$ be a random variable uniformly distributed in the interval $(0, 2\pi)$,
and define
\[
  z_i = \cos(iw) \quad (i = 1, 2, \ldots).
\]
It can be shown that $\E(z_i) = 0$, $\Var(z_i) = 1/2$, and
$\Cov(z_i, z_j) = 0$ for $i \neq j$.  So $\{z_i\}$ is white noise.  However,
clearly, it is not an independent white noise process.  It is not even strictly
stationary.
\end{example}

%% ─── Ergodicity ───────────────────────────────────────────────────────────
\paragraph{Ergodicity.}
A stationary process $\{z_i\}$ is said to be \textbf{ergodic} if, for any two bounded
functions $f\colon \mathbb{R}^k \to \mathbb{R}$ and
$g\colon \mathbb{R}^\ell \to \mathbb{R}$,
\[
  \lim_{n \to \infty}
  \big|\E[f(z_i, \ldots, z_{i+k})\,g(z_{i+n}, \ldots, z_{i+n+\ell})]\big|
  = \big|\E[f(z_i, \ldots, z_{i+k})]\big|\,
    \big|\E[g(z_{i+n}, \ldots, z_{i+n+\ell})]\big|.
\]
Heuristically, a stationary process is ergodic if it is asymptotically independent,
that is, if any two random variables positioned far apart in the sequence are almost
independently distributed.  A stationary process that is ergodic will be called
\textbf{ergodic stationary}.  Ergodic stationarity will be an integral ingredient in
developing large-sample theory because of the following property.

\medskip
\noindent\textbf{Ergodic Theorem:}\quad
\emph{(See, e.g., Theorem~9.5.5 of Karlin and Taylor (1975).) Let
$\{\mathbf{z}_i\}$ be a stationary and ergodic process with
$\E(\mathbf{z}_i) = \bm{\mu}$.\footnote{So the mean is assumed to exist and is
finite.}  Then}
\[
  \bar{\mathbf{z}}_n \equiv \frac{1}{n} \sum_{i=1}^{n} \mathbf{z}_i
  \asto \bm{\mu}.
\]

The Ergodic Theorem, therefore, is a substantial generalization of Kolmogorov's
LLN.  Serial dependence, which is ruled out by the i.i.d.\ assumption in
Kolmogorov's LLN, is allowed in the Ergodic Theorem, provided that it disappears
in the long run.  Since, for any (measurable) function $f(\cdot)$,
$\{f(\mathbf{z}_i)\}$ is ergodic stationary whenever $\{\mathbf{z}_i\}$ is, this
theorem implies that any moment of a stationary and ergodic process (if it exists
and is finite) is consistently estimated by the sample moment.  For example,
suppose $\{\mathbf{z}_i\}$ is stationary and ergodic and
$\E(\mathbf{z}_i\mathbf{z}_i')$ exists and is finite.  Then
$\frac{1}{n}\sum_i \mathbf{z}_i\mathbf{z}_i'$ is consistent for
$\E(\mathbf{z}_i\mathbf{z}_i')$.

The simplest example of ergodic stationary processes is independent white noise
processes.  (White noise processes where independence is weakened to no serial
correlation are not necessarily ergodic; Example~\ref{ex:2.4} above is an example.)
Another important example is the AR(1) process satisfying
\[
  z_i = c + \rho z_{i-1} + \varepsilon_i, \quad |\rho| < 1,
\]
where $\{\varepsilon_i\}$ is independent white noise.

%% ─── Martingales ──────────────────────────────────────────────────────────
\subsection*{Martingales}

Let $x_i$ be an element of $\mathbf{z}_i$.  The scalar process $\{x_i\}$ is called a
\textbf{martingale with respect to} $\{\mathbf{z}_i\}$ if
\begin{equation}
  \E(x_i \mid \mathbf{z}_{i-1}, \mathbf{z}_{i-2}, \ldots, \mathbf{z}_1)
  = x_{i-1} \quad \text{for } i \geq 2.\footnote{If the process started in the
  infinite past so that $i$ runs from $-\infty$ to $+\infty$, the definition is
  $\E(x_i \mid \mathbf{z}_{i-1}, \mathbf{z}_{i-2}, \ldots) = x_{i-1}$, and the
  qualifier ``$i \geq 2$'' is not needed.  Whether the process started in the infinite
  past or in date $i = 1$ is not important for large-sample theory to be developed
  below.  What will matter is that the process starts before the sample period.}
  \label{eq:2.2.4}
\end{equation}
The conditioning set $(\mathbf{z}_{i-1}, \mathbf{z}_{i-2}, \ldots, \mathbf{z}_1)$ is
often called the \textbf{information set} at point (date) $i - 1$.  $\{x_i\}$ is called
simply a \textbf{martingale} if the information set is its own past values
$(x_{i-1}, x_{i-2}, \ldots, x_1)$.  If $\mathbf{z}_i$ includes $x_i$, then $\{x_i\}$
is a martingale, because
\begin{align*}
  \E(x_i \mid x_{i-1}, x_{i-2}, \ldots, x_1)
  &= \E[\E(x_i \mid \mathbf{z}_{i-1}, \mathbf{z}_{i-2}, \ldots, \mathbf{z}_1)
     \mid x_{i-1}, x_{i-2}, \ldots, x_1] \\
  &\quad \text{(Law of Iterated Expectations)} \\
  &= \E(x_{i-1} \mid x_{i-1}, x_{i-2}, \ldots, x_1) = x_{i-1}.
\end{align*}

A vector process $\{\mathbf{z}_i\}$ is called a \textbf{martingale} if
\begin{equation}
  \E(\mathbf{z}_i \mid \mathbf{z}_{i-1}, \ldots, \mathbf{z}_1)
  = \mathbf{z}_{i-1} \quad \text{for } i \geq 2.
  \label{eq:2.2.5}
\end{equation}

\begin{example}[Hall's Martingale Hypothesis]
\label{ex:2.5}
Let $\mathbf{z}_i$ be a vector containing a set of macroeconomic variables (such as
the money supply or GDP) including aggregate consumption $c_i$ for period $i$.
Hall's (1978) martingale hypothesis is that consumption is a martingale with respect
to $\{\mathbf{z}_i\}$:
\[
  \E(c_i \mid \mathbf{z}_{i-1}, \mathbf{z}_{i-2}, \ldots, \mathbf{z}_1) = c_{i-1}.
\]
This formalizes the notion in consumption theory called ``consumption smoothing'':
the consumer, wishing to avoid fluctuations in the standard of living, adjusts
consumption in date $i-1$ to the level such that no change in subsequent consumption
is anticipated.
\end{example}

%% ─── Random Walks ─────────────────────────────────────────────────────────
\paragraph{Random Walks.}
An important example of martingales is a \textbf{random walk}.  Let $\{\mathbf{g}_i\}$
be a vector independent white noise process (so it is i.i.d.\ with mean $\mathbf{0}$
and finite variance matrix).  A random walk, $\{\mathbf{z}_i\}$, is a sequence of
cumulative sums:
\begin{equation}
  \mathbf{z}_1 = \mathbf{g}_1,\;
  \mathbf{z}_2 = \mathbf{g}_1 + \mathbf{g}_2,\; \ldots,\;
  \mathbf{z}_i = \mathbf{g}_1 + \mathbf{g}_2 + \cdots + \mathbf{g}_i, \;\ldots\;.
  \label{eq:2.2.6}
\end{equation}
Given the sequence $\{\mathbf{z}_i\}$, the underlying independent white noise
sequence, $\{\mathbf{g}_i\}$, can be backed out by taking first differences:
\begin{equation}
  \mathbf{g}_1 = \mathbf{z}_1,\;
  \mathbf{g}_2 = \mathbf{z}_2 - \mathbf{z}_1,\; \ldots,\;
  \mathbf{g}_i = \mathbf{z}_i - \mathbf{z}_{i-1},\; \ldots\;.
  \label{eq:2.2.7}
\end{equation}
So the first difference of a random walk is independent white noise.  A random walk
is a martingale because
\begin{align}
  \E(\mathbf{z}_i \mid \mathbf{z}_{i-1}, \ldots, \mathbf{z}_1)
  &= \E(\mathbf{z}_i \mid \mathbf{g}_{i-1}, \ldots, \mathbf{g}_1)
     \quad (\text{since } (\mathbf{z}_{i-1}, \ldots, \mathbf{z}_1) \text{ and}
     \notag \\
  &\qquad (\mathbf{g}_{i-1}, \ldots, \mathbf{g}_1) \text{ have the same information,
     as just seen}) \notag \\
  &= \E(\mathbf{g}_1 + \mathbf{g}_2 + \cdots + \mathbf{g}_i \mid
     \mathbf{g}_{i-1}, \ldots, \mathbf{g}_1) \notag \\
  &= \E(\mathbf{g}_i \mid \mathbf{g}_{i-1}, \ldots, \mathbf{g}_1)
     + (\mathbf{g}_1 + \cdots + \mathbf{g}_{i-1}) \notag \\
  &= \mathbf{g}_1 + \cdots + \mathbf{g}_{i-1} \quad
     (\E(\mathbf{g}_i \mid \mathbf{g}_{i-1}, \ldots, \mathbf{g}_1) = \mathbf{0}
     \text{ as } \{\mathbf{g}_i\} \text{ is independent white noise}) \notag \\
  &= \mathbf{z}_{i-1} \quad (\text{by the definition of } \mathbf{z}_{i-1}).
  \label{eq:2.2.8}
\end{align}

%% ─── Martingale Difference Sequences ──────────────────────────────────────
\subsection*{Martingale Difference Sequences}

A vector process $\{\mathbf{g}_i\}$ with $\E(\mathbf{g}_i) = \mathbf{0}$ is called a
\textbf{martingale difference sequence} (\textbf{m.d.s.}) or \textbf{martingale
differences} if the expectation conditional on its past values, too, is zero:
\begin{equation}
  \E(\mathbf{g}_i \mid \mathbf{g}_{i-1}, \mathbf{g}_{i-2}, \ldots, \mathbf{g}_1)
  = \mathbf{0} \quad \text{for} \quad i \geq 2.
  \label{eq:2.2.9}
\end{equation}
The process is so called because the cumulative sum $\{\mathbf{z}_i\}$ created from
a martingale difference sequence $\{\mathbf{g}_i\}$ is a martingale; the proof is the
same as in \eqref{eq:2.2.8}.  Conversely, if $\{\mathbf{z}_i\}$ is martingale, the
first differences created as in \eqref{eq:2.2.7} are a martingale difference sequence.

A martingale difference sequence has no serial correlation (i.e.,
$\Cov(\mathbf{g}_i, \mathbf{g}_{i-j}) = \mathbf{0}$ for all $i$ and $j \neq 0$).
A proof of this claim is as follows.

\begin{proof}
First note that we can assume, without loss of generality, that $j \geq 1$.  Since
the mean is zero, it suffices to show that
$\E(\mathbf{g}_i\mathbf{g}_{i-j}') = \mathbf{0}$.  So consider rewriting it as
follows.
\begin{align*}
  \E(\mathbf{g}_i\mathbf{g}_{i-j}')
  &= \E[\E(\mathbf{g}_i\mathbf{g}_{i-j}' \mid \mathbf{g}_{i-j})]
     \quad \text{(by the Law of Total Expectations)} \\
  &= \E[\E(\mathbf{g}_i \mid \mathbf{g}_{i-j})\,\mathbf{g}_{i-j}']
     \quad \text{(by the linearity of conditional expectations)}.
\end{align*}
Now, since $j \geq 1$, $(\mathbf{g}_{i-1}, \ldots, \mathbf{g}_{i-j}, \ldots,
\mathbf{g}_1)$ includes $\mathbf{g}_{i-j}$.  Therefore,
\begin{align*}
  \E(\mathbf{g}_i \mid \mathbf{g}_{i-j})
  &= \E[\E(\mathbf{g}_i \mid \mathbf{g}_{i-1}, \ldots, \mathbf{g}_{i-j}, \ldots,
     \mathbf{g}_1) \mid \mathbf{g}_{i-j}]
     \quad \text{(by the Law of Iterated Expectations)} \\
  &= \mathbf{0}.
\end{align*}
The last equality holds because
$\E(\mathbf{g}_i \mid \mathbf{g}_{i-1}, \ldots, \mathbf{g}_{i-j}, \ldots,
\mathbf{g}_1) = \mathbf{0}$.
\end{proof}

%% ─── ARCH Processes ───────────────────────────────────────────────────────
\subsection*{ARCH Processes}

An example of martingale differences, frequently used in analyzing asset returns,
is an autoregressive conditional heteroskedastic (\textbf{ARCH}) process introduced
by Engle (1982).  A process $\{g_i\}$ is said to be an \textbf{ARCH process of
order 1} (\textbf{ARCH(1)}) if it can be written as
\begin{equation}
  g_i = \sqrt{\zeta + \alpha g_{i-1}^2} \cdot \varepsilon_i,
  \label{eq:2.2.10}
\end{equation}
where $\{\varepsilon_i\}$ is i.i.d.\ with mean zero and unit variance.  If $g_1$ is
the initial value of the process, we can use \eqref{eq:2.2.10} to calculate
subsequent values of $g_i$.  For example,
\[
  g_2 = \sqrt{\zeta + \alpha g_1^2} \cdot \varepsilon_2.
\]
More generally, $g_i$ $(i \geq 2)$ is a function of $g_1$ and
$(\varepsilon_2, \varepsilon_3, \ldots, \varepsilon_i)$.  Therefore, $\varepsilon_i$
is independent of $(g_1, g_2, \ldots, g_{i-1})$.  It is then easy to show that
$\{g_i\}$ is an m.d.s.\ because
\begin{align}
  \E(g_i \mid g_{i-1}, g_{i-2}, \ldots, g_1)
  \label{eq:2.2.11} \\
  &= \E\!\left(\sqrt{\zeta + \alpha g_{i-1}^2} \cdot \varepsilon_i
     \mid g_{i-1}, g_{i-2}, \ldots, g_1\right) \notag \\
  &= \sqrt{\zeta + \alpha g_{i-1}^2}\;
     \E(\varepsilon_i \mid g_{i-1}, g_{i-2}, \ldots, g_1) \notag \\
  &= \sqrt{\zeta + \alpha g_{i-1}^2}\; \E(\varepsilon_i)
     \quad (\text{since } \varepsilon_i \text{ is independent of }
     (g_1, g_2, \ldots, g_{i-1})) \notag \\
  &= 0 \quad (\text{since } \E(\varepsilon_i) = 0).
  \label{eq:2.2.12}
\end{align}

By a similar argument, it follows that
\begin{equation}
  \E(g_i^2 \mid g_{i-1}, g_{i-2}, \ldots, g_1)
  = \zeta + \alpha g_{i-1}^2.
  \label{eq:2.2.13}
\end{equation}
So the conditional second moment (which equals the conditional variance since
$\E(g_i \mid g_1, g_2, \ldots, g_{i-1}) = 0$) is a function of its own history of
the process.  In this sense the process exhibits \textbf{own conditional
heteroskedasticity}.  It can be shown (see, e.g., Engle, 1982) that the process is
strictly stationary and ergodic if $|\alpha| < 1$, provided that $g_1$ is a draw
from an appropriate distribution or provided that the process started in the infinite
past.  If $g_i$ is stationary, the \emph{unconditional} second moment is easy to
obtain.  Taking the unconditional expectation of both sides of \eqref{eq:2.2.13} and
noting that
\[
  \E[\E(g_i^2 \mid g_{i-1}, g_{i-2}, \ldots, g_1)] = \E(g_i^2)
  \quad \text{and} \quad
  \E(g_i^2) = \E(g_{i-1}^2) \text{ if } g_i \text{ is stationary},
\]
we obtain
\begin{equation}
  \E(g_i^2) = \zeta + \alpha\,\E(g_i^2) \quad \text{or} \quad
  \E(g_i^2) = \frac{\zeta}{1 - \alpha}.
  \label{eq:2.2.14}
\end{equation}
If $\alpha > 0$, this model captures the characteristic found for asset returns that
large values tend to be followed by large values.  (For more details of ARCH
processes, see, e.g., Hamilton, 1994, Section~21.1.)

%% ─── Different Formulation of Lack of Serial Dependence ───────────────────
\subsection*{Different Formulation of Lack of Serial Dependence}

Evidently, an independent white noise process is a stationary martingale difference
sequence with finite variance.  And, as just seen, a martingale difference sequence
has no serial correlation.  Thus, we have three formulations of a lack of serial
dependence for zero-mean covariance stationary processes.  They are, in the order
of strength,
\begin{align}
  &(1)\;\text{``$\{\mathbf{g}_i\}$ is independent white noise.''} \notag \\
  \Rightarrow\; &(2)\;\text{``$\{\mathbf{g}_i\}$ is stationary m.d.s.\ with finite
  variance.''} \notag \\
  \Rightarrow\; &(3)\;\text{``$\{\mathbf{g}_i\}$ is white noise.''}
  \label{eq:2.2.15}
\end{align}

Condition (1) is stronger than (2) because there are processes satisfying (2) but not
(1).  An ARCH(1) process \eqref{eq:2.2.10} with $|\alpha| < 1$ is an example.
Figure~2.1 shows how a realization of a process satisfying (1) typically differs from
that satisfying (2).  Figure~2.1, Panel~(a), plots a realization of a sequence of
independent and normally distributed random variables with mean 0 and unit variance.
Panel~(b) plots an ARCH(1) process \eqref{eq:2.2.10} with $\zeta = 0.2$ and
$\alpha = 0.8$ (so that the unconditional variance, $\zeta/(1-\alpha)$, is unity as
in panel~(a)), where the value of the i.i.d.\ sequence $\varepsilon_i$ in
\eqref{eq:2.2.10} is taken from Figure~2.1, Panel~(a), so that the sign in both
panels is the same at all points.  The series in Panel~(b) is generally less volatile
than in Panel~(a), but at some points it is much more volatile.  Nevertheless, the
series is stationary.

Condition (2) is stronger than (3); the process in Example~\ref{ex:2.4} is white
noise, but (as you will show in a review question) it does not satisfy (2).

%% ── Figure 2.1 placeholder ────────────────────────────────────────────────
% Figure 2.1: Plots of Serially Uncorrelated Time Series:
%   (a) An i.i.d. N(0,1) sequence.  (b) ARCH(1) with shocks taken from panel (a).

%% ─── The CLT for Ergodic Stationary Martingale Differences Sequences ──────
\subsection*{The CLT for Ergodic Stationary Martingale Differences Sequences}

The following CLT extends the Lindeberg--Levy CLT to stationary and ergodic m.d.s.

\medskip
\noindent\textbf{Ergodic Stationary Martingale Differences CLT (Billingsley, 1961):}
\quad\emph{Let $\{\mathbf{g}_i\}$ be a vector martingale difference sequence that is
stationary and ergodic with
$\E(\mathbf{g}_i\mathbf{g}_i') = \bm{\Sigma}$,\footnote{Since $\{\mathbf{g}_i\}$ is
stationary, this matrix of cross moments does not depend on $i$.  Also, since a cross
moment matrix is indicated, it is implicitly assumed that all the cross moments exist
and are finite.} and let
$\bar{\mathbf{g}} \equiv \frac{1}{n}\sum_{i=1}^{n} \mathbf{g}_i$.  Then}
\[
  \sqrt{n}\,\bar{\mathbf{g}}
  = \frac{1}{\sqrt{n}} \sum_{i=1}^{n} \mathbf{g}_i
  \dto N(\mathbf{0},\, \bm{\Sigma}).
\]

Unlike in the Lindeberg--Levy CLT, there is no need to subtract the mean from
$\mathbf{g}_i$ because the unconditional mean of an m.d.s.\ is by definition zero.
For the same reason, $\bm{\Sigma}$ also equals $\Var(\mathbf{g}_i)$.  This CLT,
being applicable not just to i.i.d.\ sequences but also to stationary martingale
differences such as ARCH(1) processes, is more general than Lindeberg--Levy.

We have presented an LLN for serially correlated processes in the form of the
Ergodic Theorem.  A central limit theorem for serially correlated processes will be
presented in Chapter~6.
